\documentclass[12pt, titlepage]{article}
%\usepackage[norsk]{babel}
\usepackage[utf8]{inputenc}
\usepackage{minted}
\usepackage[colorlinks = true,
linkcolor = blue,
urlcolor = blue]{hyperref}


\begin{document}

\title{Mappeoppgave - del 1}
\author{IIRA2001 - Høst 2024}
\date{}
\maketitle

\section*{Mappeoppgave del 1}
Første del av mappen dreier seg om generell programmering.
Du skal legge ved tre små programmer om temaer fra faget så langt:
\begin{itemize}
  \item {Logistisk vekst}
  \item {\textit{Én av:}}
     \begin{itemize}
        \item {Sparekalkulator} 
        \item {\textit{eller} Databehandling av kunder og pengefordeling}
     \end{itemize}
  \item {\textit{Ett simuleringsprogram}}
     \begin{itemize}
        \item{Simulasjon av kundeadferd i butikk}
        \item{\textit{eller} Simulasjon av markedsdynamikk}
     \end{itemize}
\end{itemize}

Til programmene skriver dere også en rapport. Her forklarer dere hva hensikten med programmet er, eller hvilket
problem det forsøkker å løse, og hvordan dette gjøres i ulike deler av programmene.

I tillegg trenger vi en resultat- og diskusjonsdel hvor dere trekker frem eventuelle resultat fra programmet, og diskuterer disse.
Denne delen er spesielt viktig; den muliggjør at dere kan vise dere i stand til å bruke python som et verktøy i en faglig sammenheng, og 
dere får muligghet til å utvise selvstendighet, evne til kritisk tenkning og analyse. Dette er høye læringsmål som gir stor uttelling til en sensur.

Dere kan skrive rapporten i markdown celler i samme notebook som programkoden deres, eller dere kan levere rapporten inn som en egen fil ved siden av (*.pdf)

\subsection*{Kodeinnhold}

I første del av mappen ser vi etter at dere kan : 
\begin{itemize}
  \item {lage egne funksjoner for å }
     \begin{itemize}
        \item {unngå å gjenta kode}
        \item {gi programmet en god struktur}
        \item {kunne gjenbruke deler av programmet i andre sammenhenger}
     \end{itemize}
  \item {bruke variabler på en fornuftig måte}
  \item {behandle data ved bruk av lister eller dictionaries}
  \item {manipulere disse dataene med \mintinline{python}{for} eller \mintinline{python}{while} løkker}
  \item {styre programflyt med \mintinline{python}{if}/\mintinline{python}{else}}
  \item {plotte resultat med \mintinline{python}{pyplot}}
  \item {\textbf{Bruke python som verktøy i en faglig sammenheng}}
\end{itemize}

Dere \textbf{må} ikke bruke alle konseptene vi har sett på så langt, \textit{det aller viktigste er at dere kan 
bruke python til å løse et praktisk problem eller undersøke en relevant problemstilling.}

Det eneste vi forventer \textbf{skal} være med - er en fornuftig fremstilling av resultatene.
Det vil som oftest si et plot/figur av en eller form, men om dette ikke passer, kunne det for eksempel
være en tabell eller en tekst output med \mintinline{python}{print} 

\section*{Innhold}

\subsection*{Populasjonsvekst/Logistisk vekst}
Oppgaver 1.1, 2.1, og 3.1.1-3.1.4 omhandlet populasjonsvekst og/eller den logistiske formelen.
Som et minimum legger dere ved løsningen deres på en av disse oppgavene.

Dersom dere vil forbedre dette, kan dere ta i bruk programmet til å undersøke noe konkret.
Kan dere for eksempel si noe om populasjonsvekst på Sunnmøre eller i Norge?

Den logistiske likningen brukes også til å modellere hvordan ny teknologi tas i bruk (se \href{https://en.wikipedia.org/wiki/Diffusion_of_innovations}{Diffusion of innovation} og \href{https://no.wikipedia.org/wiki/Diffusjon_(kommunikasjon)}{Diffusjon (norsk wiki)})
Kanskje kan du si noe om hvor mange el-biler man kan forvente på veiene fremover?

Den logistiske modellen kan være en god kandidat for vekst som starter eksponentielt men minker etterhvert som praktiske begrensninger svekker vekstraten.
Oppgave 3 inneholder noen eksempler


\subsection*{Sparing og databehandling}

Dere skal legge ved \textbf{enten} en sparekalkulator \textbf{eller} et program som undersøker datasettet \mintinline{python}{kundedata1.json}

\subsubsection*{Sparekalkulator}

Oppgaver 1.2, 2.2 og 3.2.1-3.2.3 handlet om en sparekalkulator.
Som et minimum legger dere ved løsningen deres på en av disse oppgavene.

Som en forbedring, kan dere som foreslått i oppgaven også ta med inflasjon, formueskatt og skatt av rentene i sparekalkulatoren.
Kommer dere på mer nyttig funksjonalitet i et program som har med sparing å gjøre, kan dere implementere disse,
eller kanskje dere kan undersøke noen konkrete tilfeller som dere synes er interessante (Er du kanskje \href{https://www.imdb.com/title/tt0584425/}{Futurama} fan?)

\subsubsection*{Databehandling av kunder i sparebank}

I oppgave 4 så vi på hvordan fordelingen av sparepenger blant en kundegruppe på 2000 kunder, og på hvordan den  forandrer seg over tid.

Oppgave 4.2 bygde videre på 4.1, 4.3 på 4.2 osv, slik at det blir naturlig å levere et program som løser flere av oppgavene i Oppgave 4.

Dersom dere vil gjøre noe utover oppgavebeskrivelsene i 4 kan dere:
\begin{itemize}
   \item{Utvide datasettet}
      \begin{itemize}
         \item{Gi kundene data om yrke -- la ulike yrker ha ulik startfordeling av sparepenger/inntekt}
         \item{Gi kundene data om utdanning -- la ulike utdanningsløp eller lengde tilsvare ulike inntekt- sparebetingelser}
      \end{itemize}
   \item{Lag ekstra funksjonalitet}
      \begin{itemize}
         \item {Du kan lage funksjoner som søker opp kunder og viser relevant info og plot om de}
         \item {Lag funksjoner som lar deg hente ut og vise data om ulike kundesegment (saldo større enn 'X' osv), litt som et kundefilter}
      \end{itemize}
   \item {Plot kumulativ distribusjon og sammenlign med ulike teoretiske distribusjoner (normal, exponensial..)}
   \item {Lag funksjoner som regner ut statistiske nøkkeltall (gjennomsnitt, varians, standardavvik osv)}
\end{itemize}

Dere står fritt til å leke dere med og manipulere datasettet.
Det vi vil se i denne oppgaven er at dere kan navigere gjennom og bruke litt vanskeligere datastrukturer (dictionaries).
Det er også viktig å vise at man kan trekke ut informasjon fra datasettet, gjennom feks histogram og kakediagram.

\subsection*{Simuleringsoppgave}

Oppgave 5 dreier seg om simulering. Dere står fritt til å velge om dere vil simulere kundeadferd (5.1) eller markedsdynamikk (5.2),
men dere burde ha hatt mikroøkonomi (AE201615) om dere vil simulere markedsdynamikk

Som en minimum legger dere her ved løsning av oppgave 5.1 \textit{eller} 5.2.
Dere oppfordres til å være litt kreative og utvide simuleringsparameterene

\subsubsection*{5.1 - Kundeadferd}
Oppgave 5.1 går ut på å undersøke hvordan omsetning henger sammen med salgspris når varer selges med rabatterte priser.
Vi kan anta at man ser at man selger flere farer jo bedre rabatter man gir, men oppgaven sier ingenting om overskudd.
En naturlig forlengelse ville være å gjøre noen antagelser om dekningsgraden eller innkjøpsprisen til varene, og finne ut hvilken "rabatt"
som maksimerer overskuddet.

\subsubsection*{5.2 - Markedsdynamikk}

Oppgave 5.2 går ut på å simulere markedsdynamikk for et gode med linjære markedsmodeller for tilbud og etterspørsel.
Oppgaven slik gitt beregner ny likevektspris og kvantum for forskjellige markedssjokk - og halvveis i simuleringen innføres en intervensjon.
Som et minimum kan du løse oppgaven slik den er gitt -- en sterkt anbefalt forbedring ville være å sette simuleringen i en konkret kontekst.
Velg et gode med fornuftige verdier for parametrene $a,b,c,d$ og gjør konkrete tolkninger for betydningen av simuleringsprogrammet

Ellers kan man omtrent kaste læreboken fra mikroøkonomi (Pindyck, Rubinfeld) på simuleringen og leke seg med å undersøke mye.
Man kan feks:
\begin{itemize}
   \item {Justere $b$ og $d$ over tid for å simulere endringer i priselastisitet}
   \item {Simulere flere goder og ha med krysspriselastisitet (substitutter eller komplementariteter)}
   \item {Beregn konsumentoverskudd og/eller produsentoverskudd for hver periode}
\end{itemize}

\section*{Tilbakemelding}

Se egen info om hvordan underveisvurdering og tilbakemelding skal gjennomføres

\end{document}
